\section{Related Work}

\subsection{Deep Learning for Rice Disease Classification}

Deep learning and transfer learning have become common approaches for rice disease classification because they can learn discriminative visual patterns directly from leaf images. A DenseNet121-based model for rice plant disease classification reported an average testing accuracy of 97.9\%, showing that pre-trained CNN backbones can provide strong performance in this domain \cite{ismail2026densenet}. This result is relevant to the present study because it confirms that CNN-based models are already strong baselines for rice disease recognition. However, such works mainly evaluate the final predictive capacity of a deep architecture, rather than isolating the contribution of classical preprocessing or handcrafted descriptors under the same experimental protocol.

\subsection{Segmentation-Aware Plant Disease Recognition}

Segmentation-aware approaches attempt to reduce the influence of background and irrelevant visual regions by focusing the classifier on the leaf or lesion area. A contour-driven segmentation approach combined with optimized transfer learning architectures achieved 98.43\% test accuracy and 98.70\% F1-score for rice leaf disease classification \cite{shakeel2026contour}. This antecedent is directly related to the segmentation component of the present work, since both approaches use region-focused processing before classification. Nevertheless, the cited study evaluates segmentation as part of a larger optimized transfer learning pipeline, while this paper evaluates segmentation as an isolated ablation factor against an RGB baseline and texture-based variants.

\subsection{Explicit Texture and Hybrid Feature Descriptors}

Texture-based descriptors remain relevant in rice disease classification because symptoms such as spots, lesions, necrotic regions, and discoloration can be represented through spatial intensity relationships and local texture patterns. A feature-fusion method combining GLCM, LBP, and color histograms achieved 97.22\% accuracy on RLD1 and 99.41\% on RLD2, outperforming several individual descriptor combinations \cite{sutikno2026glcm_lbp_color}. Similarly, a hybrid feature extraction approach based on GLCM, GLDM, texture, FFT, and DWT features with a feed-forward neural network achieved an average accuracy of 94.84 $\pm$ 0.04\%, with the authors also reporting statistical validation through a paired test \cite{prity2026ffnn}. These studies support the use of explicit handcrafted descriptors for capturing complementary information, especially texture and frequency patterns that may not be explicitly controlled in standard CNN pipelines.

Overall, previous studies show that CNN backbones, segmentation-aware processing, and handcrafted texture or frequency descriptors can each improve plant or rice disease recognition. However, they are usually evaluated as independent solutions or as part of broader optimized architectures. This makes it difficult to determine the individual and combined contribution of segmentation and explicit texture features when integrated with a lightweight CNN. Therefore, this paper focuses on a controlled ablation study that compares four configurations under the same dataset split, training setup, metrics, and multi-seed evaluation protocol: RGB baseline, segmented-input CNN, CNN with texture fusion, and a segmentation-texture hybrid model.
